\section{Propiedad intelectual}

El proyecto se acoge al Reglamento de Propiedad Intelectual de la Universidad EAFIT \cite{eafit_pi}. Los \emph{derechos morales} sobre el trabajo de grado, el firmware y la documentación son perpetuos e inalienables de Alexander Saldarriaga Vélez como autor, y del asesor en lo relativo a su aporte intelectual. Los \emph{derechos patrimoniales} sobre la obra —documento, planos, código fuente y resultados— corresponden a la Universidad EAFIT en los términos del Reglamento, por tratarse de una obra desarrollada en actividad académica institucional y con su infraestructura; los materiales aportados por el estudiante no alteran esta distribución. La protección aplicable es la de \emph{derechos de autor}, sin registro constitutivo. El prototipo se transfiere al Laboratorio de Mecatrónica para uso docente e interno, sin contraprestación. No se prevé patente: la novedad está en la configuración y el ajuste experimental de tecnologías conocidas.
